% \tracingoutput SHOWS EACH PAGE AS IT IS SHIPPED.
%
%   <blank line>
%   Completed box being shipped out [1]
%   \hbox(2.0+0.0)x3.4
%   .\rule(2.0+*)x0.4
%   .\kern 3.0
%   <blank line>
%
% The heading stands on a line of its own with a blank line above it, and
% then the page's box under \showboxdepth and \showboxbreadth like any
% other box -- so a negative \showboxdepth writes
%
%   Completed box being shipped out [7.0.3] []
%
% on the one line.
%
% THE PAGE NUMBERS ARE THE SAME ONES the plain progress mark uses:
% \count0 through the last of \count0..\count9 that is not zero, joined
% by dots. `[1]', `[7.0.3]'.
%
% WHERE THE CLOSING BRACKET GOES IS THE PART THAT IS EASY TO MISS. With
% \tracingoutput the `]' is written straight after the counts, before the
% page is written out; without it, the mark opens with `[' before the
% page is written and closes with `]' after -- which is what puts the
% file names a \write opens inside the brackets. Same two characters,
% different places.
\tracingonline=1 \showboxdepth=5 \showboxbreadth=10
\message{[A tracingoutput=1]}
\tracingoutput=1
\shipout\hbox{\vrule height2pt\kern3pt}
\message{[B tracingoutput=0]}
\tracingoutput=0
\shipout\hbox{\vrule height2pt}
\message{[C counts set]}
\tracingoutput=1 \count0=7 \count1=0 \count2=3 \count9=0
\shipout\hbox{\vrule height1pt}
\message{[D showboxdepth=-1]}
\showboxdepth=-1 \showboxbreadth=-1
\shipout\hbox{\vrule height1pt\kern2pt}
\message{[E]}
\end
